\section{Results}
\label{sec:results}

\subsection{Derivative accuracy does not remove availability delay}
\label{sec:results-derivatives}

% CLAIM: C05
\Cref{fig:derivative-timing} summarizes the information required by the three
finite-difference stencils, while \cref{tab:derivative-accuracy} reports their
errors against analytic truth.  Across the three references, offline centered
velocity estimates reduce derivative RMSE by factors from
\CenteredVelocityAccuracyFactorMin{} to
\CenteredVelocityAccuracyFactorMax{} relative to the corresponding backward
differences; acceleration factors range from
\CenteredAccelerationAccuracyFactorMin{} to
\CenteredAccelerationAccuracyFactorMax{}.  Those accuracy gains do not make
the offline estimate causal: its value represented at \(t_k\) requires
\(p_{k+1}\).  The causal centered construction instead carries one sample of
native delay before propagation.  These are derivative-source properties, not
position-following lags.

\begin{table}[t]
  \centering
  \caption{Derivative accuracy on the three analytic references.  Velocity
  RMSE is in \si{\radian\per\second}; acceleration RMSE is in
  \si{\radian\per\second\squared}.  ``Centered offline'' consumes one future
  sample, whereas ``centered causal'' has one-sample native delay and is
  propagated to the latest sample.}
  \label{tab:derivative-accuracy}
  \small
  % Generated; do not edit.
\begin{tabular}{lllrr}
\toprule
Reference & Derivative source & Causal & $v$ RMSE & $a$ RMSE \\
\midrule
Quadratic & Backward & yes & 0.012136 & 0.079125 \\
Quadratic & Centered offline & no & 0.000132 & 0.000250 \\
Quadratic & Centered causal & yes & 0.000264 & 0.079125 \\
Cubic & Backward & yes & 0.004576 & 0.041940 \\
Cubic & Centered offline & no & 0.000070 & 0.000176 \\
Cubic & Centered causal & yes & 0.000140 & 0.041940 \\
Sine & Backward & yes & 0.014593 & 0.150332 \\
Sine & Centered offline & no & 0.000251 & 0.000747 \\
Sine & Centered causal & yes & 0.000501 & 0.150332 \\
\bottomrule
\end{tabular}

\end{table}

\subsection{Reliable velocity helps under the analytic information condition}
\label{sec:results-analytic}

% CLAIM: C03
On every smooth analytic reference, replacing P with PV truth reduces both
position RMSE and best grid lag
(\cref{fig:phase-a-ablation,tab:analytic-tracking}).  The RMSE reduction ranges from
\AnalyticVelocityImprovementMin\% to
\AnalyticVelocityImprovementMax\%.  P-only best lags span
\AnalyticPOnlyLagMinMS--\AnalyticPOnlyLagMaxMS{} ms, whereas PV truth records
\AnalyticPVLagMS{} ms on all three references.  The paired absolute RMSE
values, rather than the relative percentages alone, are given in
\cref{tab:analytic-tracking}.  This supports the value of reliable target
velocity for these three clean, pre-specified references and this fixed
follower protocol.  It does not show that an online position-only estimator
will supply velocity of comparable quality.

% CLAIM: C04
Adding analytic acceleration does not improve the recorded position result
over PV truth.  PV- and PVA-truth RMSE values differ by at most
\(\AnalyticPVPVAMaxRMSEDifference\,\si{\radian}\), and both have the same
\AnalyticPVLagMS{}-ms best lag on every reference.  Mechanistically, both
conditions reach the same requested endpoint position and velocity; the
different endpoint acceleration does not change the recorded position
endpoint in these runs.  This numerical non-result is not an equivalence
claim and provides no evidence about acceleration-estimator quality, internal
profile quality, control effort, or acceleration-active reference regimes.

\begin{figure}[t]
  \centering
  \includegraphics[width=\linewidth]{phase_a_ablation.pdf}
  \caption{Post-hoc visualization of the checked-in current Phase~A aggregates
  for three smooth, single-axis analytic references under
  \(\mathrm{target}[k]\mapsto\mathrm{output}[k+1]\).  Analytic PV and PVA are
  oracle-information conditions: velocity reduces P-only position error and
  lag, while acceleration adds no recorded position benefit beyond PV under
  this protocol.}
  \label{fig:phase-a-ablation}
\end{figure}

\begin{table}[t]
  \centering
  \caption{Position following for the analytic P/PV-truth/PVA-truth ablation.
  Values are deterministic outcomes on three references, not repeated-trial
  inferential statistics.  Truth derivatives are oracle information.}
  \label{tab:analytic-tracking}
  \small
  % Generated; do not edit.
\begin{tabular}{llrrr}
\toprule
Reference & Target & RMSE (rad) & Max. error (rad) & Lag (ms) \\
\midrule
Quadratic & P & 0.07567 & 0.13464 & 80 \\
Quadratic & PV truth & 0.00928 & 0.01506 & 10 \\
Quadratic & PVA truth & 0.00928 & 0.01506 & 10 \\
Cubic & P & 0.01362 & 0.02605 & 40 \\
Cubic & PV truth & 0.00345 & 0.00592 & 10 \\
Cubic & PVA truth & 0.00345 & 0.00592 & 10 \\
Sine & P & 0.04971 & 0.14553 & 70 \\
Sine & PV truth & 0.00675 & 0.01694 & 10 \\
Sine & PVA truth & 0.00675 & 0.01694 & 10 \\
\bottomrule
\end{tabular}

\end{table}

\subsection{The next-cycle oracle validates the index convention}

% CLAIM: C08
Supplying the analytic state that belongs to the next output time gives
\OracleLagMS{}-ms best grid lag on all three references, with worst RMSE
\SI{\OracleWorstRMSE}{\radian} and one-step reachability recorded for every evaluated
target.  This near-zero-error result is a protocol sanity control: it isolates
correct future-state timing from estimation and confirms the implemented
\(\mathrm{target}[k]\mapsto\mathrm{output}[k+1]\) convention on these feasible
references.  It consumes one future sample, is noncausal, and is neither a
deployable predictor nor a general tracking-error bound.

\subsection{Finite-difference pipelines do not improve the development trace}
\label{sec:results-csv}

% CLAIM: C06,N01
The P-only condition records \SI{\CSVPOnlyRMSE}{\radian} RMSE and
\CSVPOnlyLagMS{} ms best lag on the fixed-grid development trace
(\cref{fig:csv-negative,tab:csv-tracking}).  Every tested unfiltered
finite-difference-derived pipeline has larger RMSE, spanning
\SI{\CSVDifferenceRMSEMin}{\radian}--\SI{\CSVDifferenceRMSEMax}{\radian};
their best lags span
\CSVDifferenceLagMinMS--\CSVDifferenceLagMaxMS{} ms.  The observation is about
the complete pipeline, including the protocol's fixed admissibility projection
when triggered and subsequent \OrdinaryRuckig{} following.  It is one
development trace with no derivative truth, so it cannot support
generalization to independent trajectories or show that derivative
information is intrinsically unhelpful.

% CLAIM: C07
P and all PV rows require no target projection, whereas every PVA
finite-difference row projects \CSVProjectionRate\% of requested targets.  The
raw differentiated acceleration peak is
\SI{\CSVRawAccelerationPeak}{\radian\per\second\squared}, approximately
\CSVAccelerationLimitMultiple{} times the configured acceleration limit.
These values diagnose the pre-projection target estimates, not a measured
robot, simulated plant, or executed-command acceleration violation.  Because
projection is absent from P/PV but active for all PVA rows, the design cannot
separate estimator error from the projection intervention.

\begin{figure}[t]
  \centering
  \includegraphics[width=\linewidth]{csv_negative_result.pdf}
  \caption{Current negative evidence on one \CSVRowCount-row development CSV replayed
  on a fixed \SI{\ControlPeriodMS}{\milli\second} grid.  All tested unfiltered
  finite-difference-derived pipelines have higher RMSE than P.  P/PV project
  no targets, while every PVA variant projects \CSVProjectionRate\%; the
  resulting pipeline comparison cannot isolate differentiation from
  projection.  Offline-centered rows are noncausal.}
  \label{fig:csv-negative}
\end{figure}

\begin{table}[t]
  \centering
  \caption{Development-CSV position metrics and pre-following target-projection
  rates.  The CSV has no derivative truth, and ``offline'' consumes a future
  sample.  Results apply only to this fixed-grid development pipeline.}
  \label{tab:csv-tracking}
  \small
  % Generated; do not edit.
\begin{tabular}{llrrrr}
\toprule
Target & Info. & RMSE (rad) & Max. error (rad) & Lag (ms) & Proj. (\%) \\
\midrule
P & causal & 0.03519 & 0.18453 & 70 & 0.00 \\
PV backward & causal & 0.06398 & 0.32713 & 110 & 0.00 \\
PVA backward & causal & 0.03874 & 0.24049 & 70 & 32.64 \\
PV centered offline & offline & 0.06723 & 0.33287 & 130 & 0.00 \\
PVA centered offline & offline & 0.04470 & 0.26078 & 80 & 32.64 \\
PV centered causal & causal & 0.07856 & 0.46619 & 160 & 0.00 \\
PVA centered causal & causal & 0.04411 & 0.24313 & 80 & 32.64 \\
\bottomrule
\end{tabular}

\end{table}

\subsection{Frozen direct-command audit and measured runtime}
\label{sec:results-v3}

% CLAIM: C10
Within the frozen synthetic v3 protocol, the direct method records zero
continuous V/A/internal-J violations, zero projections, and zero fallbacks
over \VThreeDirectCycleCount{} command cycles
(\cref{fig:v3-direct,tab:v3-safety}).  All \VThreeDirectCycleCount{} command
states are point-admissible and one-step reachable, and all
\VThreeDirectTransitionCount{} adjacent
nonfallback transitions are sequence-consistent.  Maximum velocity and
acceleration reach their configured limits, so their minimum margins are
exactly zero rather than strictly positive; they pass the registered
\(\VThreeAuditTolerance\) audit tolerance.  Maximum internal jerk remains below the
configured jerk limit.  These counts describe the frozen direct
constant-jerk method label and do not attribute the observation to its PVA
target components.

The separately defined post-warm-up runtime population contains
\VThreeRuntimeCycleCount{} cycles (\cref{tab:v3-runtime}).  End-to-end runtime
has p99 \SI{\VThreeRuntimePNinetyNineUS}{\micro\second}, maximum
\SI{\VThreeRuntimeMaxUS}{\micro\second}, and zero recorded
\SI{\ControlPeriodMS}{\milli\second} deadline misses.  This is an empirical
distribution on one recorded software/hardware environment, not a
worst-case-execution-time statement.  Neither the safety counts nor the
runtime counts establish plant-dynamics, hardware, collision, or production
safety, and zero recorded events are not a universal failure-probability
claim.

\begin{figure}[t]
  \centering
  \includegraphics[width=\linewidth]{v3_direct_safety_runtime.pdf}
  \caption{Post-hoc visualization of frozen synthetic v3 evidence.  The
  constant-jerk command audit uses \VThreeDirectCycleCount{} cycles; the
  runtime panel uses a distinct \VThreeRuntimeCycleCount{}-cycle post-warm-up
  population.  Zero recorded violations or misses are scoped observations,
  not hardware certification or universal guarantees.}
  \label{fig:v3-direct}
\end{figure}

\begin{table}[t]
  \centering
  \caption{Frozen direct-command audit.  Sequence consistency uses
  \VThreeDirectTransitionCount{}
  adjacent transitions; other rows use \VThreeDirectCycleCount{} command
  cycles.  Continuous constraint acceptance permits margin
  \(\geq\VThreeAuditTolerance\); observed V/A minimum margins are zero.}
  \label{tab:v3-safety}
  \small
  % Generated; do not edit.
\begin{tabular}{lrrl}
\toprule
Frozen v3 audit quantity & Observed & Denominator & Status \\
\midrule
Continuous V/A/internal-J violations & 0.000 & 42199 & pass \\
Point-admissible rate & 1.000 & 42199 & pass \\
One-step-reachable rate & 1.000 & 42199 & pass \\
Sequence-consistent rate & 1.000 & 42079 & pass \\
Projected cycles & 0.000 & 42199 & pass \\
\bottomrule
\end{tabular}

\end{table}

\begin{table}[t]
  \centering
  \caption{Frozen direct-method runtime after per-trajectory warm-up.  This
  population is separate from the command-profile audit and reports measured,
  not worst-case, execution time.}
  \label{tab:v3-runtime}
  \small
  % Generated; do not edit.
\begin{tabular}{lrrl}
\toprule
Frozen v3 runtime quantity & Observed & Denominator & Population \\
\midrule
Timed cycles & 30199 & -- & locked-test, after warm-up \\
Total compute p99 ($\mu$s) & 874.92 & 30199 & locked-test, after warm-up \\
Total compute maximum ($\mu$s) & 2113.25 & 30199 & locked-test, after warm-up \\
10-ms deadline misses & 0 & 30199 & locked-test, after warm-up \\
\bottomrule
\end{tabular}

\end{table}

\subsection{Profile-aware auditing changes what can be claimed}

% CLAIM: C11,C12
Current profile-aware infrastructure represents an exposed native
\Ruckig{} control-period prefix as ordered piecewise-constant-jerk segments
and integrates those segments individually.  If an exact internal segment
model is unavailable, the record marks internal jerk unavailable rather than
substituting average acceleration-difference jerk.  The same records separate
native \OrdinaryRuckig{}, \ViabilityShieldedRuckig{}, and
\DirectConstantJerkExecution{}, including any algorithm-changing fallback.
This makes method semantics auditable in current code.  It is not a v3
rerun, a corrected comparative outcome, or evidence that one follower
outperforms another.
